\documentclass[10pt]{article}

% ---------- Document Identity (update these only) ----------
\newcommand{\DocStage}{}
\newcommand{\DocTitle}{Your Road to Tier One SOF}
\newcommand{\ProgramName}{182-Week Civilian-to-Selection Readiness Pipeline}
\newcommand{\ProgramSubtitle}{}

% ---------- Page ----------
\usepackage[legalpaper,left=0.5in,right=0.5in,top=0.6in,bottom=0.45in]{geometry}

% ---------- Typography ----------
\usepackage[T1]{fontenc}
\usepackage[utf8]{inputenc}
\usepackage{libertinus}
\usepackage{microtype}
\usepackage{parskip}
\usepackage{ragged2e}
\usepackage{textcomp}
\AtBeginDocument{\color{black}}
\AtBeginDocument{\pagecolor{white}}
\AtBeginDocument{\justifying}

% ---------- Landscape pages ----------
\usepackage{pdflscape}

% ---------- Color ----------
\usepackage[table]{xcolor}
\definecolor{StageBlue}{HTML}{1F4E79}
\definecolor{StageBlueDark}{HTML}{163A5A}
\definecolor{LightGray}{HTML}{F2F4F7}
\definecolor{MidGray}{HTML}{D9DEE5}
\definecolor{RuleGray}{HTML}{C7CED8}
\definecolor{HeaderInk}{HTML}{000000}
\definecolor{Ink}{HTML}{000000}

% ---------- Utility ----------
\usepackage{enumitem}
\sloppy
\setlength{\emergencystretch}{2em}

% ---------- Boxes / UI ----------
\usepackage[most]{tcolorbox}

\tcbset{
  charterTitle/.style={
    enhanced,
    colback=StageBlue!4,
    colframe=StageBlue,
    boxrule=1.25pt,
    arc=3pt,
    left=16pt,right=16pt,top=14pt,bottom=14pt,
    borderline north={3.2pt}{0pt}{StageBlue},
    borderline west={4pt}{0pt}{StageBlue},
    borderline south={1.25pt}{0pt}{StageBlue},
    drop shadow={black!25!white},
  },
  charterRule/.style={
    enhanced,
    colback=LightGray,
    colframe=RuleGray,
    boxrule=0.85pt,
    arc=3pt,
    left=12pt,right=12pt,top=10pt,bottom=10pt,
    borderline west={3pt}{0pt}{StageBlue},
  },
  charterBadge/.style={
    enhanced,
    colback=StageBlue,
    coltext=white,
    colframe=StageBlueDark,
    boxrule=0.6pt,
    arc=2pt,
    left=6pt,right=6pt,top=3pt,bottom=3pt,
  }
}
\newtcbox{\Badge}{nobeforeafter,charterBadge}

% ---------- Lists ----------
\setlist[itemize]{leftmargin=1.5em, itemsep=0.25em, topsep=0.25em}
\setlist[enumerate]{leftmargin=1.8em, itemsep=0.25em, topsep=0.25em}

% ---------- TOC ----------
\usepackage{tocloft}
\setcounter{tocdepth}{3}
\setlength{\cftbeforesecskip}{3pt}
\setlength{\cftbeforesubsecskip}{2pt}
\setlength{\cftbeforesubsubsecskip}{0pt}
\renewcommand{\cftsecfont}{\bfseries}
\renewcommand{\cftsecpagefont}{\bfseries}
\renewcommand{\cftsubsecfont}{\normalfont}
\renewcommand{\cftsubsecpagefont}{\normalfont}
\renewcommand{\cftsubsubsecfont}{\normalfont}
\renewcommand{\cftsubsubsecpagefont}{\normalfont}
\setlength{\cftsecindent}{0pt}
\setlength{\cftsubsecindent}{1.25em}
\setlength{\cftsubsubsecindent}{2.8em}
\setlength{\cftsecnumwidth}{2.1em}
\setlength{\cftsubsecnumwidth}{2.8em}
\setlength{\cftsubsubsecnumwidth}{3.6em}

% Remove the built-in TOC heading so only the custom heading is shown
\makeatletter
\renewcommand{\tableofcontents}{\@starttoc{toc}}
\makeatother

% ---------- Header/Footer: page number only ----------
\usepackage{fancyhdr}
\pagestyle{fancy}
\fancyhf{}
\renewcommand{\headrulewidth}{0pt}
\renewcommand{\footrulewidth}{0pt}
\fancyfoot[C]{\thepage}

% ---------- Section formatting ----------
\usepackage{titlesec}

\titleformat{\section}
  {\Large\bfseries\color{Ink}}
  {\thesection}{0.6em}{}
  [\vspace{0.25em}\textcolor{StageBlue}{\rule{\linewidth}{0.8pt}}]

\titleformat{\subsection}
  {\large\bfseries\color{Ink}}
  {\thesubsection}{0.6em}{}

\titleformat{\subsubsection}
  {\normalsize\bfseries\color{Ink}}
  {\thesubsubsection}{0.6em}{}

\titlespacing*{\section}{0pt}{0.95em}{0.35em}
\titlespacing*{\subsection}{0pt}{0.65em}{0.25em}
\titlespacing*{\subsubsection}{0pt}{0.55em}{0.2em}

% ---------- Table engine ----------
\usepackage{tabularray}
\UseTblrLibrary{booktabs}

% ---------- tabularray: GLOBAL table treatment ----------
\SetTblrInner{
  cells = {fg=black, bg=white, valign=t},
}

% ---------- Header helpers ----------
\newcommand{\Hdr}[1]{\shortstack[c]{\strut #1\strut}}

% ---------- Lines ----------
\newcommand{\TitleLine}{\textcolor{StageBlue}{\rule{0.98\linewidth}{0.95pt}}}
\newcommand{\SubTitleLine}{\textcolor{RuleGray}{\rule{0.98\linewidth}{0.65pt}}}

% ---------- Continued on next page ----------
\newcommand{\ContinuedOnNextPageInline}{%
  \par
  \vfill
  \noindent\textcolor{RuleGray}{\rule{\linewidth}{1pt}}\vspace{-2pt}
  \noindent\hfill\textit{Continued on next page}\par\vspace{-10pt}
  \noindent\textcolor{RuleGray}{\rule{\linewidth}{1pt}}
  \clearpage
}

% ---------- PDF hyperlinks ----------
\usepackage[hidelinks]{hyperref}
\hypersetup{
  colorlinks=false,
  pdfborder={0 0 0},
  linktoc=all,
  pdftitle={\DocTitle},
  pdfauthor={\ProgramName}
}

% =========================================================
\begin{document}

\subsection*{Phase 05 Card Header}
\phantomsection
\addcontentsline{toc}{subsection}{Phase 05 Card Header}

\begin{itemize}
\item Phase \textbf{ID}: Phase 05
\item Phase Name: Aerobic Base and Extended Stamina Development
\item Duration weeks: 12
\item Version: v1.0
\item Stage Alignment: Stage 5
\item Governing Status: Deployable
\end{itemize}

\subsection*{1. Phase Mission}
\phantomsection

Build extended aerobic stamina and repeatable endurance verification discipline while maintaining calisthenics continuity and conservative load tolerance, producing gate-grade evidence only inside protected deload and test windows.

\subsection*{2. Phase Purpose}
\phantomsection

\begin{itemize}
\item Establish longer-duration aerobic output capability with low variance and without validity drift from fatigue contamination.
\item Introduce the first consequential extended \textbf{RUN} verification event (5-mile) as an end-phase gate-grade construct while preserving safety and reslot discipline.
\item Maintain \textbf{CAL} battery continuity and maintain \textbf{RUCK} tolerance progression as supporting gate context without stacking high-cost domains unsafely.
\item Set up Phase 06 by improving endurance repeatability, recovery stability, and comparability discipline for longer tests.
\end{itemize}

\subsection*{3. Entry Criteria}
\phantomsection

\begin{itemize}
\item Entry requires admissible evidence from Phase 04 exit posture (\textbf{VALID} tests and admissible session compliance) per Stage 1.F.
\item Durability posture must be stable enough to tolerate increasing time-on-feet exposure:
\begin{itemize}
\item no unresolved gait-altering pain posture,
\item foot hotspot/blister management is stable and logged.
\end{itemize}
\item If evidence is incomplete or invalid, default posture is \textbf{HOLD} and reslot per Stage 3; do not “start anyway” on missing gate evidence.
\end{itemize}

\subsection*{4. Operating Constraints}
\phantomsection

\begin{itemize}
\item Six sessions per week are mandatory.
\item Required session elements per Stage 1.F in correct order:
\begin{itemize}
\item Safety self-check first
\item Warm-up
\item Mobility
\item Main work
\item Cardio
\item Cooldown
\item Recovery notes and any required post-session notes per Stage 1.F
\end{itemize}
\item Cardio minimum standard: minimum 30 minutes continuous per session.
\item 10,000 steps per day is a global requirement and must be logged.
\item Validity doctrine: admissible \textbf{VALID} evidence only for gates; invalid tests provide no gate credit.
\item Phase 05 does not alter Stage 3 protected windows; gate-grade tests occur only in deload and test windows.
\item Water governance and safety boundaries remain controlling; \textbf{SWIM} tests occur only when pool verification and safety posture permit valid execution.
\end{itemize}

\subsection*{5. Primary Adaptations and Physiological Focus}
\phantomsection

\begin{itemize}
\item Extended aerobic stamina: improved tolerance for longer steady outputs with stable gait mechanics.
\item Efficiency and pacing control under non-breathless posture in build weeks; gate-grade tests are protected from fatigue contamination.
\item Foot and lower-limb tissue tolerance: stable foot interface under increasing time-on-feet exposure.
\item Recovery stability: reduced variance in sleep/fatigue markers so longer tests can be executed as \textbf{VALID} when scheduled.
\item \textbf{CAL} continuity: maintain calisthenics capacity without allowing \textbf{CAL} fatigue to contaminate extended \textbf{RUN} verification windows.
\end{itemize}

\subsection*{6. Primary Modalities}
\phantomsection

\begin{itemize}
\item Emphasized domains: \textbf{RUN}; \textbf{CAL}; \textbf{RECOVERY}; \textbf{DURABILITY}; \textbf{BODYCOMP}.
\item Supporting domains: \textbf{RUCK} (supporting tolerance and periodic verification per phase map); \textbf{SWIM} (periodic verification only if pool access is stable and validity can be preserved).
\item De-emphasized/prohibited posture:
\begin{itemize}
\item no new tests or protocols,
\item no stacking high-cost major tests beyond Stage 3.A doctrine intent,
\item no unsupervised aquatic exposures; underwater is not part of Phase 05 gate posture.
\end{itemize}
\end{itemize}

\clearpage
\begin{landscape}

\subsection*{7. Weekly Structure}
\phantomsection

\begingroup
\scriptsize
\setlength{\tabcolsep}{2pt}
\renewcommand{\arraystretch}{1.12}

\begin{longtblr}{
  width=\linewidth,
  rowhead=1,
  colspec = {
    Q[l,wd=0.70in]
    X[l,1.60]
    X[l,1.10]
    X[l,2.70]
    X[l,2.45]
  },
  hlines,
  vlines,
  cells = {hsep=6pt, vsep=6pt, halign=l, valign=t},
  cell{1-Z}{1-5} = {fg=black},
  row{odd} = {bg=white},
  row{even} = {bg=LightGray},
  row{1} = {bg=StageBlue!15, fg=black, font=\bfseries, halign=l, valign=t},
}
\Hdr{Session \#} &
\Hdr{Primary Focus} &
\Hdr{Main Work Domains} &
\Hdr{Cardio Modality/Intent} &
\Hdr{Notes/Risk Controls} \\

1 &
Aerobic base throughput &
\textbf{RUN}; \textbf{DURABILITY} &
Modality: step-based locomotion default; Minutes: 30–60; RPE plus Talk Test: RPE 3–5, full-to-short sentences; Continuity: continuous planned &
Low-impact bias; protect gait stability; downshift if foot hotspots trend worse. \\

2 &
Calisthenics continuity &
\textbf{CAL}; \textbf{DURABILITY} &
Modality: low-impact steady; Minutes: 30–45; RPE plus Talk Test: RPE 2–4, full sentences; Continuity: continuous planned &
Treat \textbf{CAL} as controlled exposure; avoid breathless density that contaminates later endurance work. \\

3 &
Endurance support day &
\textbf{RUN}; \textbf{RECOVERY} &
Modality: step-based or treadmill as controlled substitute; Minutes: 30–60; RPE plus Talk Test: RPE 3–5, full-to-short sentences; Continuity: continuous planned &
No novelty near protected windows (footwear, route, or tool changes). \\

4 &
Durability and movement quality maintenance &
\textbf{DURABILITY}; \textbf{CAL} &
Modality: lowest-risk steady; Minutes: 30–45; RPE plus Talk Test: RPE 2–4, full sentences; Continuity: continuous planned &
Downshift-first posture if sleep/fatigue trends degrade; preserve continuous Cardio. \\

5 &
Load tolerance supporting exposure &
\textbf{RUCK}; \textbf{DURABILITY} &
Modality: low-impact steady; Minutes: 30–60; RPE plus Talk Test: RPE 3–5, full-to-short sentences; Continuity: continuous planned &
\textbf{RUCK} remains controlled; avoid stacking with the longest \textbf{RUN} verification inside the same protected window. \\

6 &
Consolidation and pre-window stabilization &
\textbf{RECOVERY}; \textbf{DURABILITY} &
Modality: step-based default; Minutes: 30–45; RPE plus Talk Test: RPE 2–4, full sentences; Continuity: continuous planned &
Enforce no-novelty posture approaching Week 4/8/12 protected windows; preserve validity readiness. \\

\end{longtblr}

\endgroup

\subsection*{8. Progression Rules}
\phantomsection

\begin{itemize}
\item One progression lever at a time: do not increase more than one primary lever per week across longest Cardio-primary duration, total weekly Cardio minutes (session minutes only), and complexity/density of main work.
\item Protect high-cost verification: build weeks emphasize stability; do not create “informal time trials” outside protected windows.
\item Regression triggers: durability drift (hotspots, gait changes), recovery collapse (sleep streak, rising fatigue), environment hazards, or tool drift trigger downshift or substitution rather than frequency reduction.
\item Substitution posture: preserve intent, reduce risk, and keep evidence loggable; substitutions that increase risk are prohibited.
\item Invalid evidence is non-admissible for gates; reslot per Stage 3 rather than stacking make-ups.
\end{itemize}

\subsection*{9. Deload and Test Placement}
\phantomsection

\begin{itemize}
\item Protected deload and test windows (week-in-phase): Week 4 (mid-phase touchpoint); Week 8 (mid-phase touchpoint); Week 12 (end-phase exit gate).
\item Gate-grade tests are executed only inside these protected windows.
\item Staging posture: if multiple domains are assessed in a protected week, apply staged battery intent consistent with Stage 3 anti-stacking/spacing; prioritize the dominant gate item and reslot secondary tests if spacing cannot be protected.
\end{itemize}

\end{landscape}

\clearpage
\begin{landscape}

\subsection*{10. Test Battery}
\phantomsection

\begingroup
\scriptsize
\setlength{\tabcolsep}{2pt}
\renewcommand{\arraystretch}{1.12}

\begin{longtblr}{
  width=\linewidth,
  rowhead=1,
  colspec = {
    X[l,1.20]
    X[l,2.45]
    Q[l,wd=0.85in]
    X[l,2.65]
    Q[l,wd=1.20in]
  },
  hlines,
  vlines,
  cells = {hsep=6pt, vsep=6pt, halign=l, valign=t},
  cell{1-Z}{1-5} = {fg=black},
  row{odd} = {bg=white},
  row{even} = {bg=LightGray},
  row{1} = {bg=StageBlue!15, fg=black, font=\bfseries, halign=l, valign=t},
}
\Hdr{Test Timing} &
\Hdr{Test \textbf{ID}/Name} &
\Hdr{Domain} &
\Hdr{Validity Conditions} &
\Hdr{Pass Threshold Stage 2 Tier} \\

Week 4 Mid-phase &
\textbf{C10} / \textbf{MET-2B-016} 3-Mile Run Time Trial &
\textbf{RUN} &
Stage 2 protocol validity; route or treadmill admissibility requires full environment/tool logging; no stop-rule violations &
\textbf{INTERMEDIATE} \\

Week 4 Mid-phase &
\textbf{C4} / \textbf{MET-2B-010} Push-Ups — 2-Minute Timed, Strict &
\textbf{CAL} &
Stage 2 \textbf{CAL} artifact evidence required for gate-validity; missing artifact = \textbf{INVALID} for gate use &
\textbf{INTERMEDIATE} \\

Week 4 Mid-phase &
\textbf{C5} / \textbf{MET-2B-011} Sit-Ups — 2-Minute Timed, Standard &
\textbf{CAL} &
Stage 2 \textbf{CAL} artifact evidence required for gate-validity; missing artifact = \textbf{INVALID} for gate use &
\textbf{INTERMEDIATE} \\

Week 4 Mid-phase &
\textbf{C6} / \textbf{MET-2B-012} Pull-Ups — Strict Dead-Hang &
\textbf{CAL} &
Stage 2 \textbf{CAL} artifact evidence required for gate-validity; missing artifact = \textbf{INVALID} for gate use &
\textbf{INTERMEDIATE} \\

Week 4 Mid-phase &
\textbf{C7} / \textbf{MET-2B-013} Plank — Forearm &
\textbf{CAL} &
Stage 2 \textbf{CAL} artifact evidence required for gate-validity; missing artifact = \textbf{INVALID} for gate use &
\textbf{INTERMEDIATE} \\

Week 8 Mid-phase &
\textbf{C12} / \textbf{MET-2B-019} 6-Mile Ruck Time Trial — 35 lb dry &
\textbf{RUCK} &
Dry load verified; water logged separately; environment/tool logging complete; no stop-rule violations &
\textbf{INTERMEDIATE} \\

Week 8 Mid-phase &
\textbf{C13} / \textbf{MET-2B-022} Swim 500 yd Time Trial — No Fins &
\textbf{SWIM} &
Pool unit/length verified; no fins; allowable strokes; swim governance rules honored; no stop-rule violations &
\textbf{INTERMEDIATE} \\

Week 12 End-phase &
\textbf{C11} / \textbf{MET-2B-017} 5-Mile Run Time Trial &
\textbf{RUN} &
Route/tool comparability protected; no novelty near gate; environment/tool logging complete; no stop-rule violations &
\textbf{INTERMEDIATE} \\

Week 12 End-phase &
\textbf{C4} / \textbf{MET-2B-010} Push-Ups — 2-Minute Timed, Strict &
\textbf{CAL} &
Stage 2 \textbf{CAL} artifact evidence required; missing artifact = \textbf{INVALID} for gate use &
\textbf{INTERMEDIATE} \\

\end{longtblr}

\endgroup

\end{landscape}
\clearpage

\subsection*{11. Exit Standards}
\phantomsection

\begin{itemize}
\item Phase 05 exit gate posture aligns to \textbf{INTERMEDIATE} tier theme.
\item \textbf{ADVANCE} requires admissible evidence only:
\begin{itemize}
\item End-phase gate tests in Week 12 are \textbf{VALID} and meet \textbf{INTERMEDIATE} tier threshold posture per Stage 2.
\item Any required gate test that is \textbf{INVALID} defaults the metric posture to \textbf{HOLD} until valid retest exists; reslot per Stage 3.
\end{itemize}
\item Phase 05 exit decision is constrained by Stage 1.F decision-window compliance and durability posture:
\begin{itemize}
\item unresolved gait-altering pain posture, active blister, or repeated stop-rule events constrain \textbf{ADVANCE}.
\end{itemize}
\end{itemize}

\subsection*{12. Risk Controls and Stop Rules}
\phantomsection

\begin{itemize}
\item Stage 0.5 and Stage 1.F remain controlling for stop rules and escalation posture; do not continue a test or session through red-flag symptoms.
\item Phase 05 high-risk controls:
\begin{itemize}
\item extended \textbf{RUN} verification (\textbf{C11} / \textbf{MET-2B-017}) is protected from stacking with the longest/heaviest \textbf{RUCK} in the same protected week; apply staged battery intent and reslot secondary tests if needed.
\item Foot integrity is a primary constraint: hotspot-to-blister progression triggers downshift and reslot posture; do not force gait-compromised gate tests.
\item Environmental hazards (ice, unsafe heat, poor air quality) trigger postponement or controlled substitution only when comparability and logging preserve validity; otherwise reslot.
\item \textbf{SWIM} tests are executed only when pool verification is stable; if pool access/conditions compromise validity, reslot rather than forcing a compromised attempt.
\end{itemize}
\end{itemize}

\clearpage
\begin{landscape}

\subsection*{13. Required Capabilities and Dependencies}
\phantomsection

\begingroup
\scriptsize
\setlength{\tabcolsep}{2pt}
\renewcommand{\arraystretch}{1.12}

\begin{longtblr}{
  width=\linewidth,
  rowhead=1,
  colspec = {
    X[l,1.35]
    X[l,2.35]
    Q[l,wd=1.25in]
    Q[l,wd=1.10in]
    X[l,2.95]
  },
  hlines,
  vlines,
  cells = {hsep=6pt, vsep=6pt, halign=l, valign=t},
  cell{1-Z}{1-5} = {fg=black},
  row{odd} = {bg=white},
  row{even} = {bg=LightGray},
  row{1} = {bg=StageBlue!15, fg=black, font=\bfseries, halign=l, valign=t},
}
\Hdr{Dependency \textbf{ID}} &
\Hdr{Required Capability Category and Tier} &
\Hdr{Due-By week-in-phase} &
\Hdr{Owner} &
\Hdr{Fail Action} \\

\textbf{DEP-1C-007} \& \textbf{DEP-1C-008} &
7.11 Footwear Systems and Foot Care — Tier required for extended stamina verification &
Week 1 &
Trainee &
\textbf{HOLD} high-cost \textbf{RUN}/\textbf{RUCK} gates until foot integrity is stable; reslot tests if hotspots/blisters emerge; document validity impact. \\

\textbf{DEP-1C-006} &
7.18 Testing, Timing, Measurement, and Course-Setup Tools — Tier required for \textbf{RUN}/\textbf{RUCK} validity &
Week 1 &
Coach Lead &
If route/tool comparability cannot be preserved, tests are \textbf{INVALID} for gate use; reslot to next protected window. \\

\textbf{DEP-1C-013} &
7.07 Load-Bearing Rucking and Carry Systems — Tier required for \textbf{RUCK} \textbf{TT} validity &
Week 4 &
Equipment Lead &
If load verification or safe configuration is not possible, do not attempt gate; reslot; no proxy equivalency. \\

\textbf{DEP-1C-011} &
7.17 Aquatic Training and Water Safety Equipment — Tier required for \textbf{SWIM} \textbf{TT} validity &
Week 4 &
Coach Lead &
If pool unit and length cannot be verified or facility constraints compromise validity, do not attempt; reslot; no substitutions treated as \textbf{SWIM} evidence. \\

\textbf{DEP-1C-017} &
7.12 Recovery, Mobility, and Tissue-Care Implements — Tier supporting durability stability &
Week 1 &
Trainee &
If recovery drift prevents valid gate attempts, downshift and reslot; do not compress or stack make-ups. \\

\textbf{DEP-1C-018} &
7.13 Nutrition Preparation and Hydration Systems — Tier supporting safe extended efforts &
Week 1 &
Trainee &
If heat risk or dehydration risk is present, downshift intensity and reslot tests if safety is ambiguous; document constraints. \\

\end{longtblr}

\endgroup

\subsection*{14. Validity Requirements}
\phantomsection

\begin{itemize}
\item Admissibility is governed by Stage 1.F and Stage 2 validity rules.
\item Tests must be tagged \textbf{VALID} or \textbf{INVALID}; \textbf{INVALID} provides no gate credit and is treated as not yet tested for gate decisions.
\item \textbf{CAL} gate-validity requires compliant artifact evidence per Stage 2 protocols; missing artifact = \textbf{INVALID} for gate use.
\item Environment/tool identity and logging completeness are mandatory for \textbf{RUN}/\textbf{RUCK}/\textbf{SWIM} comparability; missing required fields defaults to \textbf{INVALID} for gate use.
\item Gate decisions use admissible evidence only inside protected windows; build-week attempts are not treated as gate-grade by default.
\end{itemize}

\end{landscape}

\subsection*{15. Change Control Notes}
\phantomsection

\begin{itemize}
\item Allowed without patch:
\begin{itemize}
\item reslotting invalid/missed tests to the next protected window within the phase per Stage 3.E,
\item swapping the order of tests within a protected week while preserving caps/spacing intent and validity posture.
\end{itemize}
\item Requires patch:
\begin{itemize}
\item adding a test week,
\item extending phase duration,
\item changing gate focus tests,
\item changing work:deload pattern.
\end{itemize}
\item Any patch is governed by Stage 1.F change control and must be crosswalk-reviewed in Stage 8.
\end{itemize}

\subsection*{16. Compliance Checklist}
\phantomsection

\begin{itemize}
\item Phase \textbf{ID}, name, duration, and governing status are stated and correct for Phase 05.
\item Gates reference Stage 2 tests and metrics only and are scheduled only in protected deload and test windows.
\item Weekly structure table is complete and uses Cardio cell structure (Modality; Minutes; RPE plus Talk Test; Continuity continuous planned).
\item Operating constraints state Cardio minimum standard and 10,000 steps per day in body text.
\item Dependencies use Stage 7 category IDs and do not require 7.02; owners are from the locked taxonomy.
\item Governing Status is Deployable because dependency IDs are explicit Stage 1 dependency identifiers in gate-relevant fields.
\end{itemize}

\end{document}